\section{F-SM-32: the interval theorem with exact coproducts}

\subsection{Statement and remaining historical scope}
Let \(k\) be a commutative ring.  Let \(\mathcal M\) be a
\(k\)-linear abelian category with enough projectives and all small
coproducts.  Suppose that
\(\varphi:\mathcal M\rightleftarrows\operatorname{Mod}k:\kappa\)
is an adjunction, with \(\kappa\) exact and coproduct preserving.
Write \(\mathcal D_{\mathcal M}=D_{\ge0}(\mathcal M)_\infty\),
\(\mathcal D_k=D_{\ge0}(k)_\infty\), and
\[
 F(X)=\mathbb L\varphi(X[0]),\qquad
 W=\{f:F(f)\text{ is an equivalence}\}.
\]
Assume:
\begin{enumerate}
\item the \emph{canonical} derived counit
\(F(\kappa k)\to k[0]\) is an equivalence;
\item each \(X\in\mathcal M\) admits a monomorphism
\(a_X:X\hookrightarrow A_X\) with \(F(A_X)=0\), including
degree-zero homology;
\item for every small family \((X_i)\), the canonical map
\[
       \bigoplus_i F(X_i)\longrightarrow
                         F\!\left(\bigoplus_iX_i\right)
                         \quad\text{is an equivalence}.       \tag{SM32.1}
\]
\end{enumerate}
Then \(F\) induces an equivalence
\[
        \mathcal M[W^{-1}]_\infty\simeq\mathcal D_k.            \tag{SM32.2}
\]
Consequently its homotopy category is the ordinary localization
\(W^{-1}\mathcal M\simeq D_{\ge0}(k)\).

In particular, condition (SM32.1) holds if \(\mathcal M\)
satisfies \emph{AB4}: its small coproduct functors are exact.
Thus the full interval theorem holds for every such AB4 category
with enough projectives, without AB5, a generator, local
presentability, an exact chain-valued realization, or a
functorial choice of projective resolutions.

The original interval supplies the second assumption by
\(A_X=J\otimes X\) and \(a_X=\alpha\otimes1_X\).
Fresh visual reading of cote~134-8, footers~46--48, confirms
the scalar adjoints, exactness of \(\kappa\), and the bilinear
tensor convention.  These pages do not explicitly require
exactness of arbitrary copowers on \emph{every} object of
\(\mathcal M\), or state AB4.  Bilinearity and the scalar
adjoints alone are not treated as such a requirement here.
The result therefore removes the AB5/presentability gate and
reduces the remaining abstract case to the precise coproduct
condition (SM32.1); it does not silently add AB4 to the original
handwritten conditions (a)--(d).

Fix universes \(\mathbb U\in\mathbb V\), with the categories
\(\mathbb V\)-small and the specified coproducts
\(\mathbb U\)-small.  Localizations are initially formed in
\(\mathbb V\).  The proof does not assume that their mapping
spaces are \(\mathbb U\)-small; the final equivalence supplies
that conclusion when \(\operatorname{Mod}k\) is formed in
\(\mathbb U\).  Geometric realizations have the countable
indexing category \(\Delta^{\mathrm{op}}\).

\subsection{Known inputs}
We use the following primary results, with their hypotheses
kept distinct from the argument below.
\begin{itemize}
\item Cisinski, \emph{Higher categories and homotopical algebra},
author version dated July~30, 2026, Proposition~7.5.6,
printed pp.~372--373, and its dual: the localization of a
category of cofibrant objects has finite colimits, and a
functor preserving its initial object and cofibration pushouts
induces a finite-colimit-preserving functor.  Proposition~7.7.1,
p.~396, handles products of such relative categories.
\item Ramzi, Sosnilo, and Winges (2025), Definition~2.5 and
Theorem~2.6(1),(4), pp.~5--6: in the same finite cofibration
setting, coproducts stable under weak equivalences pass to the
Dwyer--Kan localization.  The proof of Theorem~2.6, together
with Remark~2.3 and Example~2.4, gives the left-fraction
description used below.  Neither a sequential-colimit axiom
nor accessibility is a hypothesis here.
\item Lurie, \emph{Higher algebra} (September~18, 2017),
Theorem~1.3.3.8, p.~101: for an abelian category with enough
projectives, restriction to its projectives identifies functors
from its connective derived \(\infty\)-category preserving
geometric realizations with arbitrary functors on those
projectives.  This is an equivalence of \emph{functor
\(\infty\)-categories}, including natural transformations.
Lemma~1.3.3.11(2), pp.~103--104, and Proposition~1.3.3.16,
pp.~106--107, show that connective derived functors coming
from right \(t\)-exact exact functors preserve these
realizations.  Lemma~1.3.3.17 and its proof, p.~107, identify
the realization of the Dold--Kan simplicial object of a
connective projective complex with that complex, including
the augmentation comparison.  Theorem~1.3.4.4, p.~109,
identifies the ordinary chain-category localization.
\item Lurie, \emph{Higher topos theory},
Propositions~4.4.2.6--7, pp.~296--297: small coproducts and
pushouts generate small colimits, including the corresponding
preservation statement.
\end{itemize}
These are existing background theorems.  The complete deduction
of (SM32.2), especially its two coherent inverse comparisons,
is given next.

\subsection{Complete proof}

\paragraph{1. Exact coproducts imply (SM32.1).}
Suppose AB4.  For each \(i\), choose a projective resolution
\(P_{i,\bullet}\to X_i\).  A coproduct of projectives is
projective: maps from each summand lift separately through an
epimorphism.  AB4 makes \(\bigoplus_iP_{i,\bullet}\) a
projective resolution of \(\bigoplus_iX_i\).
As a left adjoint, \(\varphi\) preserves coproducts, so applying
it to this resolution gives
\[
 \varphi\!\left(\bigoplus_iP_{i,\bullet}\right)
                  =\bigoplus_i\varphi(P_{i,\bullet}).
\]
Coproducts in \(\mathcal D_k\) are computed by these ordinary
coproducts of complexes.  The comparison maps induced by the
summand inclusions are precisely those in (SM32.1).  This proves
the claimed sufficient condition, without any filtered-colimit
argument.  The remainder of the proof only uses (SM32.1).

\paragraph{2. A finite cofibration structure and its localization.}
Take every monomorphism of \(\mathcal M\) to be a cofibration.
Initial maps are cofibrations, and their pushouts exist and are
monomorphisms.  The coherent derived functor
\(\mathbb L\varphi\) sends a short exact sequence to its
derived cofiber sequence.  In particular, for a monomorphism
\(u\),
\[
       u\in W\quad\Longleftrightarrow\quad
                       F(\operatorname{coker}u)=0.             \tag{SM32.3}
\]
A pushout retains this cokernel, so trivial cofibrations are
stable under arbitrary pushout.  More generally, pushing a
weak equivalence along a monomorphism compares two short
exact sequences with the same quotient; their derived
triangles show that the resulting map is again in \(W\).

Every map \(f:X\to Y\) has the factorization
\[
 X\xrightarrow{(a_X,f)}A_X\oplus Y
                          \xrightarrow{\operatorname{pr}_Y}Y. \tag{SM32.4}
\]
The first map is monic, and the second is in \(W\), because
\(F\) is additive and \(F(A_X)=0\).  Choices in this
factorization need not be functorial.  Two-out-of-three and
two-out-of-six follow from the definition of \(W\).
These verifications give all the cofibrant-object axioms in
Ramzi--Sosnilo--Winges Definition~2.5.

Write \(\gamma:\mathcal M\to\mathcal C=
\mathcal M[W^{-1}]_\infty\).  The cited localization theorem
gives finite colimits in \(\mathcal C\).  By (SM32.1), small
coproducts of maps in \(W\) stay in \(W\); Theorem~2.6(4)
therefore gives small coproducts in \(\mathcal C\), preserved
by \(\gamma\).  HTT Proposition~4.4.2.6 gives all small
colimits in \(\mathcal C\).

The actual coherent functor
\(F=\mathbb L\varphi\circ(X\mapsto X[0])\) factors as
\[
             \overline F:\mathcal C\longrightarrow\mathcal D_k.
                                                               \tag{SM32.5}
\]
It preserves finite colimits.  Indeed a pushout along a
monomorphism \(A\to B\), with other map \(A\to D\), has
the exact sequence
\[
 0\longrightarrow A\longrightarrow B\oplus D
                   \longrightarrow B\amalg_A D\longrightarrow0,
\]
so \(F\) sends this pushout to a derived pushout.  Apply the
last assertion of the dual of Cisinski Proposition~7.5.6.
Equation (SM32.1) and preservation of coproducts by
\(\gamma\) show that \(\overline F\) preserves coproducts.
HTT Proposition~4.4.2.7 now shows that it preserves all small
colimits, in particular geometric realizations.

We will also need \emph{conservativity}.  The fraction
description in the cited cofibration theorem represents every
morphism of \(h\mathcal C\) by
\[
                  X\xrightarrow{f}Z\xleftarrow{s}Y,
                       \qquad s\in W.                        \tag{SM32.6}
\]
If its image under \(\overline F\) is invertible, then
\(F(f)\) is invertible because \(F(s)\) is.  Thus
\(f\in W\), and (SM32.6) is already invertible in
\(h\mathcal C\).  An \(\infty\)-morphism is an equivalence
exactly when it is invertible in its homotopy category.
Consequently \(\overline F\) is conservative.  This argument
does not infer conservativity merely from the definition of
\(W\); it uses the fraction theorem.

\paragraph{3. Scalar normalization and projective units.}
Exactness of \(\kappa\) gives, for projective \(P\in\mathcal M\),
\[
 \operatorname{Hom}_k(\varphi P,-)
                  =\operatorname{Hom}_{\mathcal M}(P,\kappa(-))
\]
exact.  Thus \(\varphi P\) is projective over \(k\).
There is the canonical derived adjunction
\(\mathbb L\varphi\dashv\kappa:\mathcal D_k\to
\mathcal D_{\mathcal M}\), with \(\kappa\) applied termwise;
its counit on a complex is represented by
\(\varphi P(\kappa C)\to\varphi\kappa C\to C\).
This notation denotes a projective representative, not a
functorial choice of one.

For a free scalar module \(k^{(I)}\), (SM32.1), preservation
of coproducts by \(\kappa\), and the normalization on \(k\)
show that the canonical counit
\[
          F(\kappa k^{(I)})\longrightarrow k^{(I)}[0]
\]
is an equivalence.  Naturality passes this assertion to every
projective scalar retract \(Q\).  These identifications are
the actual counit identifications.  Both derived functors
\(\mathbb L\varphi\kappa\) and the identity on
\(\mathcal D_k\) preserve geometric realizations, by HA
Lemma~1.3.3.11(2).  Theorem~1.3.3.8 therefore promotes the
projective identifications to the equivalence of the canonical
counit on all of \(\mathcal D_k\).

In particular, for projective \(P\in\mathcal M\), the original
adjunction unit satisfies
\[
                 \eta_P:P\longrightarrow\kappa\varphi P
                            \quad\in W.                       \tag{SM32.7}
\]
Its image under \(F\), followed by the just-established
counit identification, is the identity of \(\varphi P[0]\)
by the adjunction triangle identity.

\paragraph{4. Coherent extension from projectives.}
By HA Theorem~1.3.3.8, the functor
\(Q\mapsto\gamma\kappa Q\) on projective scalar modules
extends to a geometric-realization-preserving functor
\[
                         L:\mathcal D_k\longrightarrow\mathcal C.
                                                               \tag{SM32.8}
\]
The same theorem, applied to \(P\mapsto\gamma P\), gives
\[
                    T:\mathcal D_{\mathcal M}\longrightarrow\mathcal C.
                                                               \tag{SM32.9}
\]
All comparisons here are comparisons of \(\infty\)-functors.
Since \(\overline F\) preserves realizations, the functors
\(\overline F L\) and \(\mathrm{id}_{\mathcal D_k}\)
preserve them and have the canonically equivalent restrictions
to projective scalar modules supplied in Step~3.
Full faithfulness in HA Theorem~1.3.3.8 gives
\[
                      \overline F L\simeq\mathrm{id}_{\mathcal D_k}.
                                                               \tag{SM32.10}
\]
Likewise, \(T\) and \(L\mathbb L\varphi\) preserve
realizations.  Their restrictions to projective \(P\) are
\(\gamma P\) and \(\gamma\kappa\varphi P\), and (SM32.7)
gives a natural equivalence between them.  Thus
\[
                             T\simeq L\mathbb L\varphi.       \tag{SM32.11}
\]

For clarity, the following step constructs the required natural
augmentation; an objectwise choice of resolutions would not
by itself supply this comparison.

\paragraph{5. The comma-category augmentation and its identification.}
For a functor \(g:\mathcal A\to\mathcal E\), where both
categories admit geometric realizations, let
\((\mathrm{id}_{\mathcal E}\downarrow g)\) have objects
\((e,a,u:e\to g(a))\).  This comma \(\infty\)-category has
geometric realizations computed in its two components, even
if \(g\) does not preserve them.  Indeed a simplicial diagram
has colimit arrow
\[
 \operatorname*{colim}e_n\longrightarrow
 \operatorname*{colim}g(a_n)\longrightarrow
                                g(\operatorname*{colim}a_n).
                                                               \tag{SM32.12}
\]
To check its universal property, use the mapping-space formula
\[
 \operatorname{Map}((e,a,u),(e',a',u'))
 =\operatorname{Map}_{\mathcal E}(e,e')
  \mathop{\times}_{\operatorname{Map}_{\mathcal E}(e,g(a'))}
                 \operatorname{Map}_{\mathcal A}(a,a').        \tag{SM32.13}
\]
Mapping out of each componentwise colimit is a limit, and
limits of spaces commute with this pullback.  Hence (SM32.12)
satisfies the asserted universal property.  It also follows
that a functor into this comma category preserves realizations
whenever both its component functors do.

Apply this with \(g=\gamma:\mathcal M\to\mathcal C\),
viewing \(\mathcal M\) as its nerve.  An ordinary abelian
category has geometric realizations, computed by its reflexive
coequalizers.  Extend
\[
       P\longmapsto(\gamma P,P,\mathrm{id}_{\gamma P})
                                                               \tag{SM32.14}
\]
from the projectives by HA Theorem~1.3.3.8, now with this
comma category as target.  Its first component is \(T\).
Its second component is \(H_0:\mathcal D_{\mathcal M}
\to\mathcal M\): this functor is left adjoint to the heart
inclusion \(\iota:X\mapsto X[0]\), preserves realizations,
and agrees with the second component on projectives.
We obtain a coherent natural transformation
\[
                             b:T\longrightarrow\gamma H_0.  \tag{SM32.15}
\]

We identify its image under \(\overline F\), rather than
assuming that it is the desired augmentation.  The functor
\[
 (\mathrm{id}_{\mathcal C}\downarrow\gamma)
       \longrightarrow(\mathrm{id}_{\mathcal D_k}\downarrow F)
\]
induced by \(\overline F\) preserves realizations by
(SM32.12).  On the other hand, the functor
\[
 C\longmapsto
 \bigl(\mathbb L\varphi(C),H_0(C),
       \mathbb L\varphi(C\to\iota H_0(C))\bigr)               \tag{SM32.16}
\]
also preserves realizations: its components
\(\mathbb L\varphi\) and \(H_0\) do, and (SM32.13)
checks the comma-category claim.  Both resulting functors from
\(\mathcal D_{\mathcal M}\) to this last comma category
restrict on projective \(P\) to
\((\varphi P[0],P,\mathrm{id})\).  The full faithfulness
part of HA Theorem~1.3.3.8 identifies them as arrow-valued
functors.  In particular, when \(C=X[0]\),
\(\overline F(b_{X[0]})\) identifies with the image of the
identity truncation \(X[0]\to H_0(X[0])[0]\), and is an
equivalence.  Conservativity from Step~2 proves the natural
equivalence
\[
                             T\iota\xrightarrow{\sim}\gamma. \tag{SM32.17}
\]

Equations (SM32.11) and (SM32.17) now give a natural equivalence
\(L\overline F\gamma\simeq\gamma\).  The universal property
of relative \(\infty\)-localization is fully faithful on
functor \(\infty\)-categories; it therefore descends this
comparison to
\[
                         L\overline F\simeq\mathrm{id}_{\mathcal C}.
                                                               \tag{SM32.18}
\]
Together with (SM32.10), this proves (SM32.2).
Passing to homotopy categories gives the asserted ordinary
localization equivalence. \(\square\)

\subsection{Concrete realization and exact limitations}
If \(Q_\bullet\) is a connective complex of projective scalar
modules representing \(C\in\mathcal D_k\), the preceding
inverse is concretely
\[
               L(C)\simeq
                  \operatorname*{colim}_{\Delta^{\mathrm{op}}}
                                      \gamma\kappa\Gamma(Q)_n.
                                                               \tag{SM32.19}
\]
For a projective resolution \(P_\bullet\to X\),
\(T(X[0])\simeq\operatorname*{colim}\gamma\Gamma(P)_n\),
and (SM32.15) is its augmentation to \(\gamma X\).
These identifications follow from HA Lemma~1.3.3.17; the
comma-category argument proves that their choices fit together
coherently.  No ordinary filtered colimit in \(\mathcal M\)
is substituted for either realization.

The new scope strictly includes the AB4, non-AB5,
nonpresentable example F-SM-31.  What remains in the bare
abstract archive claim is whether its actual remaining
assumptions force (SM32.1), or whether a full tensor-and-interval
counterexample violates it.  A merely additive three-sort
construction is insufficient for that negative conclusion.
The theorem also makes no claim that an ordinary right adjoint
\(G\) to a separately chosen exact realization is already a
quasi-inverse.  Its raw counit test remains distinct.

\subsection{Primary references}
\begin{itemize}
\item Cisinski, D.-C. (2019). \emph{Higher categories and
homotopical algebra}. Cambridge University Press.
Author version dated July~30, 2026:
\url{https://cisinski.app.uni-regensburg.de/CatLR.pdf}.
Exact locators used here refer to that version:
Proposition~7.5.6, pp.~372--373; Proposition~7.7.1, p.~396;
Definition~7.7.2 and Proposition~7.7.4, p.~397.
\item Grothendieck, A. (n.d.). \emph{Cote 134-8}
[Unpublished manuscript], footers~46--48.
Archives Grothendieck, Universit\'e de Montpellier.
\url{https://grothendieck.umontpellier.fr/134-8.pdf}.
\item Lurie, J. (2009). \emph{Higher topos theory}
(Annals of Mathematics Studies, Vol.~170).
Princeton University Press.
\url{https://www.math.ias.edu/~lurie/papers/HTT.pdf}.
\item Lurie, J. (2017, September 18). \emph{Higher algebra}
[Author manuscript].
\url{https://people.math.harvard.edu/~lurie/papers/HA.pdf}.
Exact inputs: Theorem~1.3.3.8, p.~101;
Lemma~1.3.3.11(2), pp.~103--104;
Proposition~1.3.3.16 and Lemma~1.3.3.17, pp.~106--107;
Theorem~1.3.4.4, p.~109.
\item Ramzi, M., Sosnilo, V., \& Winges, C. (2025).
\emph{Every motive is the motive of a stable
\(\infty\)-category} (arXiv:2503.11338v1).
\url{https://arxiv.org/abs/2503.11338}.
Exact input: Definition~2.5 and Theorem~2.6(1),(4),
pp.~5--6, with Remark~2.3 and Example~2.4, p.~4.
\end{itemize}
